\section{Results}
\label{sec:results}

Numbers are macro-F1 (\%) averaged over a dataset's tasks unless stated otherwise. Per-task and per-dataset tables
and additional figures are in the supplementary material.

\begin{table*}[t]
\centering
\small
\caption{Clean-data benchmark across datasets: task-averaged macro-F1 (\%) at each dataset's primary label granularity. Parentheses give the selected encoder(s) and head. Rank: position of the best primary-pair head among all fusion configurations.}
\label{tab:cross_bench}
\begin{adjustbox}{max width=\textwidth}
\begin{tabular}{lcccccccc}
\toprule
Dataset & Tasks & Majority & Char TF-IDF & Best text & Best image & XLM-R+CLIP-B/32 & Best fusion & Rank \\
\midrule
CMBan & 5 & 35.0 & 46.2 & 46.9 (mE5 LR) & 51.0 (CLIP-L/14 LR) & 48.4 (Concat-LR) & 50.1 (LaBSE+CLIP-L/14 Concat-MLP) & 21/51 \\
BanglaAbuseMeme & 4 & 34.8 & 65.2 & 60.5 (XLM-R-L MLP) & 66.0 (CLIP-L/14 MLP) & 65.7 (Late) & 67.6 (mE5+CLIP-L/14 Concat-MLP) & 22/51 \\
MemoSen & 1 & 25.6 & 52.6 & 48.4 (XLM-R-L LR) & 54.4 (CLIP-L/14 MLP) & 53.6 (Concat-LR) & 56.2 (XLM-R-L+CLIP-L/14 GMU) & 16/51 \\
\bottomrule
\end{tabular}
\end{adjustbox}
\end{table*}


\subsection{Clean-Data Benchmark}
\label{sec:results:benchmark}

\textbf{Image encoders are stronger than text encoders} (Table~\ref{tab:cross_bench}). The best image model, always
CLIP ViT-L/14, beats the best text model by 4.1 (CMBan), 5.5 (BAM) and 6.0 (MemoSen) points, and CLIP-B/32 alone beats
every text encoder. Cartoon and meme templates recur, and CLIP partly reads rendered text~\cite{goh2021multimodal}.
The gap persists on MemoSen's human captions, so it is not only OCR noise.

\textbf{Frozen text encoders lose to character $n$-grams} (Table~\ref{tab:conf_text}). Character TF-IDF beats all
six text encoders on BAM (65.2 vs.\ $\leq$60.5) and MemoSen (52.6 vs.\ $\leq$48.4). Meme text is short, colloquial
and misspelt, which character $n$-grams capture directly. XLM-R-large is the best text encoder on the larger datasets
and mE5 on CMBan, where all encoders lie within 2.2 points; mBERT is the weakest. Accuracy is misleading on these
imbalanced labels: the majority class reaches 59--63\% accuracy at 26--35 macro-F1.

\begin{table}[t]
\centering
\small
\caption{Frozen text encoders: task-averaged accuracy and macro-F1 (\%) with the better head per encoder ($^{L}$: LR, $^{M}$: MLP). Bold: best text encoder. Both heads are listed in the supplementary material.}
\label{tab:conf_text}
\begin{adjustbox}{max width=\columnwidth}
\begin{tabular}{lcccccc}
\toprule
Encoder & \multicolumn{2}{c}{CMBan} & \multicolumn{2}{c}{BAM} & \multicolumn{2}{c}{MemoSen} \\
\cmidrule(lr){2-3} \cmidrule(lr){4-5} \cmidrule(lr){6-7}
 & Acc. & F1 & Acc. & F1 & Acc. & F1 \\
\midrule
XLM-R base & 50.0 & 45.9$^{L}$ & 61.5 & 59.6$^{M}$ & 57.3 & 46.0$^{L}$ \\
XLM-R large & 47.7 & 46.0$^{M}$ & 62.5 & \textbf{60.5}$^{M}$ & 60.5 & \textbf{48.4}$^{L}$ \\
mBERT & 46.0 & 44.8$^{M}$ & 59.0 & 57.2$^{M}$ & 51.7 & 41.2$^{L}$ \\
BanglaBERT & 50.7 & 46.7$^{L}$ & 60.8 & 58.5$^{L}$ & 53.4 & 42.8$^{L}$ \\
LaBSE & 46.1 & 44.7$^{M}$ & 60.2 & 58.7$^{M}$ & 54.2 & 43.3$^{L}$ \\
mE5 base & 51.9 & \textbf{46.9}$^{L}$ & 61.9 & 60.2$^{M}$ & 56.4 & 44.9$^{L}$ \\
\midrule
Majority class & 59.4 & 35.0 & 60.7 & 34.8 & 62.5 & 25.6 \\
Char TF-IDF & 53.7 & 46.2 & 68.2 & 65.2 & 67.6 & 52.6 \\
\bottomrule
\end{tabular}
\end{adjustbox}
\end{table}


\textbf{Fusion helps when enough data is available.} The best fusion model beats the best text model on every BAM
and MemoSen task (+6.9 and +7.6 points; significant after Holm correction on all five tasks) and the best image model
by +1.5 and +1.9 points (significant on two of five tasks). On CMBan no difference is significant: with 133 memes the
confidence intervals span about $\pm$8 points. The primary XLM-R + CLIP-B/32 pair ranks 16th--22nd of 51 fusion
configurations, 1.7--2.6 points behind the best, almost entirely because of its smaller image encoder.

\begin{table*}[t]
\centering
\small
\caption{Relative Robustness Drop (\%) across datasets, averaged over tasks and severities, for the matching unimodal heads (uni.) and the five XLM-R + CLIP-B/32 fusion heads (fus.; flip rate \% in parentheses). The top six rows are linguistic, the bottom six visual perturbations.}
\label{tab:cross_rob}
\begin{adjustbox}{max width=\textwidth}
\begin{tabular}{lcccccc}
\toprule
Perturbation & CMBan uni. & CMBan fus. & BAM uni. & BAM fus. & MemoSen uni. & MemoSen fus. \\
\midrule
Transliteration & 12.9 & 14.5 (37) & 18.4 & 8.2 (22) & 17.3 & 6.7 (21) \\
Roman. spelling noise & 17.9 & 20.7 (45) & 26.5 & 11.5 (28) & 21.8 & 7.7 (26) \\
Bangla spelling noise & 0.6 & 0.3 (7) & 1.7 & 0.8 (6) & 0.3 & -0.5 (8) \\
Code-mixing & -1.5 & 0.7 (8) & 0.7 & 0.3 (4) & 3.0 & 1.1 (8) \\
Emoji insertion & 2.3 & 2.0 (12) & 4.8 & 2.5 (11) & 11.7 & 4.8 (14) \\
Emoji/symbol removal & 1.0 & 0.4 (3) & 0.2 & 0.1 (1) & 3.5 & 0.9 (3) \\
Blur & 2.8 & 2.6 (23) & 16.6 & 6.3 (20) & 14.9 & 7.4 (27) \\
Crop & -1.9 & -0.7 (12) & 0.9 & 0.6 (10) & 3.3 & 1.6 (13) \\
Down-sampling & 2.2 & 1.9 (18) & 8.3 & 3.3 (15) & 13.4 & 7.1 (25) \\
JPEG & -0.0 & -0.5 (6) & 3.2 & 1.5 (8) & 0.8 & -0.1 (7) \\
Text occlusion & -0.3 & 0.3 (14) & 4.4 & 2.0 (11) & 4.0 & 1.0 (17) \\
Region occlusion & 1.4 & 0.5 (10) & 3.8 & 2.1 (11) & 4.5 & 1.6 (12) \\
\bottomrule
\end{tabular}
\end{adjustbox}
\end{table*}


\subsection{Robustness}
\label{sec:results:rob}

\textbf{Romanisation is the most damaging perturbation} on every dataset (Table~\ref{tab:cross_rob}). Romanised
spelling noise costs the fusion heads 20.7\% (CMBan), 11.5\% (BAM) and 7.7\% (MemoSen) of relative macro-F1 and
flips 26--45\% of predictions; plain transliteration costs 6.7--14.5\%. Native-script spelling noise costs at most
1.7\% and code-mixing at most 3\%, so the damage comes from the \emph{change of script}, which moves XLM-R's embedding
far from the Bangla-script training distribution. Bangla-script memes flip most under transliteration (23--43\%),
code-mixed memes less (18--28\%), and Latin-script memes almost never ($\leq$2.5\%).

\textbf{Fusion cushions text fragility only with enough data.} On BAM and MemoSen fusion halves the Romanisation
damage of the text-only heads (11.5 vs.\ 26.5\%; 7.7 vs.\ 21.8\%). On CMBan the fusion heads lose \emph{more}
(20.7 vs.\ 17.9\%): with 133 training memes they do not learn to discount unreliable text.

\textbf{Aggregate F1 hides instability.} Visual perturbations barely affect CMBan ($\leq$2.6\%), but blur and
down-sampling cost the fusion heads 3--7\% on the larger datasets, where the image models are stronger. Even small F1
changes hide many flips: blur changes 20--27\% of fusion predictions on all datasets, because broken correct
predictions are offset by repaired wrong ones. Robustness evaluations should therefore report flip rates.

\begin{table*}[t]
\centering
\small
\caption{Train-time Romanisation augmentation (XLM-R + CLIP-B/32). Macro-F1 (\%) on clean test folds and under transliteration and Romanised spelling noise (mean over heads, tasks and severities); relative Robustness Drop (\%) in parentheses. Spelling noise is unseen during training in the \emph{+translit.} setting.}
\label{tab:cross_aug}
\begin{adjustbox}{max width=\textwidth}
\begin{tabular}{lccccccc}
\toprule
Dataset & Augmentation & Text clean & Text translit. & Text noise & Fusion clean & Fusion translit. & Fusion noise \\
\midrule
CMBan & none & 45.5 & 39.7 (13) & 37.5 (18) & 47.2 & 40.6 (15) & 37.8 (21) \\
 & +translit. & 45.6 & 44.1 (3) & 43.2 (5) & 46.5 & 45.3 (2) & 44.0 (5) \\
 & +translit.\,+\,noise & 45.2 & 44.9 (1) & 45.7 (-1) & 46.1 & 46.2 (-1) & 46.0 (-0) \\
BAM & none & 59.2 & 48.2 (18) & 43.4 (26) & 65.0 & 59.8 (8) & 57.7 (12) \\
 & +translit. & 59.5 & 56.0 (6) & 53.7 (10) & 64.9 & 63.5 (2) & 63.1 (3) \\
 & +translit.\,+\,noise & 59.2 & 56.8 (4) & 55.5 (7) & 64.6 & 63.7 (2) & 63.4 (2) \\
MemoSen & none & 46.0 & 38.0 (17) & 36.0 (22) & 53.2 & 49.7 (7) & 49.1 (8) \\
 & +translit. & 45.5 & 43.2 (5) & 40.6 (11) & 53.6 & 52.4 (2) & 51.2 (4) \\
 & +translit.\,+\,noise & 45.2 & 43.7 (3) & 42.5 (6) & 53.5 & 52.7 (2) & 51.9 (3) \\
\bottomrule
\end{tabular}
\end{adjustbox}
\end{table*}


\subsection{Mitigation by Romanisation Augmentation}
\label{sec:results:aug}

Adding one transliterated copy of each training meme cuts the fusion drop under transliteration from 15/8/7\% to 2\%
on CMBan/BAM/MemoSen (Table~\ref{tab:cross_aug}). It also generalises to Romanised spelling noise, which is unseen in
this setting: the drop falls from 21/12/8\% to 5/3/4\%. Adding the noise style as well closes the remaining gap only
marginally. Clean macro-F1 changes by $-0.7$, $-0.1$ and $+0.4$ points, and the other linguistic perturbations by at
most 1.3 points, so the remedy is almost free.

\begin{table}[t]
\centering
\small
\caption{Explainability. IG text share of Concat-MLP on clean inputs and at severity 3. Ratio: Grad-CAM mass on OCR text boxes over their area share ($>1$: over-attention to rendered text). $r$: correlation of clean and perturbed CAMs (severity 2).}
\label{tab:conf_xai}
\begin{adjustbox}{max width=\columnwidth}
\begin{tabular}{lccccccc}
\toprule
 & \multicolumn{4}{c}{IG text share} & \multicolumn{3}{c}{Grad-CAM} \\
\cmidrule(lr){2-5} \cmidrule(lr){6-8}
Dataset & Clean & Translit. & Blur & Text occ. & Ratio & $r$ blur & $r$ JPEG \\
\midrule
CMBan & 0.59 & 0.73 & 0.50 & 0.51 & 1.13 & 0.30 & 0.92 \\
BAM & 0.55 & 0.71 & 0.43 & 0.44 & 0.76 & 0.39 & 0.87 \\
MemoSen & 0.55 & 0.68 & 0.46 & 0.47 & 0.88 & 0.36 & 0.92 \\
\bottomrule
\end{tabular}
\end{adjustbox}
\end{table}


\subsection{Explainability}
\label{sec:results:xai}

\textbf{Perturbed text gains attribution} (Table~\ref{tab:conf_xai}). On clean memes the text embedding receives
55--59\% of the IG attribution. Romanisation \emph{raises} this share to 0.68--0.73: the unfamiliar text vector lies
far from the IG baseline and dominates the decision instead of being discounted, which explains the large drops.
Blur and text occlusion lower the share to 0.43--0.51 by almost the same amount, suggesting that much of what the
image branch contributes is \emph{read} text.

\textbf{CLIP over-attends to text only on cartoon memes.} The Grad-CAM text-attention ratio is 1.13 on CMBan, whose
cartoon panels carry large speech bubbles, but 0.76 (BAM) and 0.88 (MemoSen), where photographs with short caption
bands dominate. The maps are stable under JPEG ($r\geq0.87$) but not under blur ($r\leq0.39$), even when the
prediction is unchanged.

\begin{table}[t]
\centering
\small
\caption{Error taxonomy of the best learned XLM-R + CLIP-B/32 fusion head (\% of memes, mean over tasks). Clean errors: contextual/cultural, linguistic, visual, OCR-limited, cross-modal. Under perturbation (severity $\geq2$): stable memes and memes broken only by linguistic, only by visual, or by both (X-mod.) perturbation types.}
\label{tab:conf_err}
\begin{adjustbox}{max width=\columnwidth}
\begin{tabular}{lccccccccc}
\toprule
 & \multicolumn{5}{c}{Clean errors} & \multicolumn{4}{c}{Under perturbation} \\
\cmidrule(lr){2-6} \cmidrule(lr){7-10}
Dataset & Ctx. & Ling. & Vis. & OCR & X-mod. & Stable & Ling. & Vis. & X-mod. \\
\midrule
CMBan & 17.6 & 8.7 & 8.9 & 6.0 & 3.6 & 11.1 & 18.5 & 6.6 & 20.5 \\
BAM & 10.8 & 7.3 & 6.3 & 6.8 & 0.9 & 32.7 & 11.1 & 10.5 & 13.6 \\
MemoSen & 10.4 & 5.7 & 5.8 & 6.3 & 3.2 & 22.2 & 14.7 & 12.7 & 19.3 \\
\bottomrule
\end{tabular}
\end{adjustbox}
\end{table}


\subsection{Error Analysis}
\label{sec:results:errors}

The largest clean error class on every dataset is \emph{contextual/cultural} (Table~\ref{tab:conf_err}): 18\% of
CMBan and 10--11\% of BAM and MemoSen memes are missed by all models although readable text is available. It peaks on
sentiment and, on CMBan, on sarcasm and humour, where memes depend on films, public figures or drifting template
meanings. Cross-modal errors are rare (1--4\%): learned fusion seldom destroys information either modality has, but
also rarely recovers memes both miss. Only 11\% (CMBan) to 33\% (BAM) of memes stay correct under every
perturbation, and memes broken by both perturbation types are the most common fragile group (14--21\%). On CMBan
Offensive, linguistic fragility reaches 44\%, as Romanisation disguises the few swear words that signal the label.
